\chapter{Guia de Implantação, DevOps e Runbook Operacional}
\label{chap:implantacao_devops}

\section{Arquitetura da Esteira de CI/CD}
A esteira de Integração e Entrega Contínua (CI/CD) automatiza o fluxo de validação e implantação a partir do repositório GitHub via \textbf{GitHub Actions}.

\begin{enumerate}
    \item \textbf{Estágio de Linting e Formatação:} Execução do Ruff (Python) e ESLint (TypeScript) para garantia de conformidade de código;
    \item \textbf{Estágio de Testes Automatizados:} Execução paralela de testes unitários e de integração com cobertura mínima exigida;
    \item \textbf{Estágio de Build e Contêinerização:} Construção das imagens Docker otimizadas multi-stage;
    \item \textbf{Estágio de Deploy Automático:} Implantação em ambiente de homologação após aprovação de Pull Request na branch principal.
\end{enumerate}

% PLACEHOLDER STARUML v7.0: Pipeline de CI/CD
% Ferramenta: StarUML v7.0 -> Model -> Add Diagram -> Activity Diagram
% Arquivo: Imagens/devops_pipeline_cicd.png (Exportar em PNG 300 DPI ou PDF vetorial)
% Descomente a linha abaixo apos salvar o arquivo no diretorio Imagens/:
% \incluirdiagrama{Imagens/devops_pipeline_cicd.png}{Fluxo Automatizado da Esteira CI/CD}{fig:devops_pipeline}
\begin{figure}[htbp]
    \centering
    \fbox{\parbox{0.9\textwidth}{\centering\vspace{1cm}%
        \textbf{[Pipeline de CI/CD -- Diagrama de Atividades -- StarUML v7.0]}\\%
        \textit{Modelar no StarUML v7.0 e exportar para: \texttt{Imagens/devops\_pipeline\_cicd.png}}\\%
        \small Menu: \texttt{File -> Export Diagram as -> PNG...} (Fundo branco, alta resolução 300 DPI)\\%
        \vspace{0.3cm}%
        \footnotesize Para ativar a imagem real, descomente no arquivo \texttt{.tex}:\\
        \texttt{\textbackslash incluirdiagrama\{Imagens/devops\_pipeline\_cicd.png\}\{Pipeline de CI/CD\}\{fig:devops_pipeline\}}%
    \vspace{1cm}}}
    \caption{Fluxo de Estágios da Esteira de CI/CD (Placeholder StarUML v7.0)}
    \label{fig:devops_pipeline}
\end{figure}

\section{Matriz de Ambientes e Gerenciamento de Configurações}
A \autoref{fig:devops_topologia} apresenta a topologia física e de execução distribuída entre os ambientes.

% PLACEHOLDER STARUML v7.0: Topologia de Ambientes e Infraestrutura
% Ferramenta: StarUML v7.0 -> Model -> Add Diagram -> Deployment Diagram
% Arquivo: Imagens/devops_topologia_infra.png (Exportar em PNG 300 DPI ou PDF vetorial)
% Descomente a linha abaixo apos salvar o arquivo no diretorio Imagens/:
% \incluirdiagrama{Imagens/devops_topologia_infra.png}{Topologia de Infraestrutura e Ambientes}{fig:devops_topologia}
\begin{figure}[htbp]
    \centering
    \fbox{\parbox{0.9\textwidth}{\centering\vspace{1cm}%
        \textbf{[Topologia de Ambientes e Infraestrutura -- Deployment Diagram -- StarUML v7.0]}\\%
        \textit{Modelar no StarUML v7.0 e exportar para: \texttt{Imagens/devops\_topologia\_infra.png}}\\%
        \small Menu: \texttt{File -> Export Diagram as -> PNG...} (Fundo branco, alta resolução 300 DPI)\\%
        \vspace{0.3cm}%
        \footnotesize Para ativar a imagem real, descomente no arquivo \texttt{.tex}:\\
        \texttt{\textbackslash incluirdiagrama\{Imagens/devops\_topologia\_infra.png\}\{Topologia de Infraestrutura\}\{fig:devops_topologia\}}%
    \vspace{1cm}}}
    \caption{Topologia de Ambientes e Infraestrutura (Placeholder StarUML v7.0)}
    \label{fig:devops_topologia}
\end{figure}

\begin{table}[htbp]
\caption{Matriz de Ambientes de Execução}
\label{tab:matriz_ambientes}
\centering
\small
\begin{tabularx}{\textwidth}{@{} l Y Y Y @{}}
\toprule
\textbf{Ambiente} & \textbf{Finalidade} & \textbf{Infraestrutura / Host} & \textbf{Estratégia de Deploy} \\
\midrule
Desenvolvimento & Testes locais dos desenvolvedores & Docker Compose local & Manual \\
Homologação & Validação de PO e testes integrados & Servidor Cloud (Staging) & Automático (branch main) \\
Produção & Ambiente operacional final em tempo real & Instâncias Cloud de Alta Disponibilidade & Manual com aprovação de 2 revisores \\
\bottomrule
\end{tabularx}
\end{table}

\section{Procedimento de Migração de Banco de Dados}
As alterações de esquema no banco PostgreSQL são gerenciadas com versionamento estrito via \textbf{Alembic}. Nunca são executadas alterações manuais diretamente no banco de produção.

\begin{lstlisting}[language=bash, caption={Comandos para Migração e Rollback de Banco de Dados}]
# Aplicar todas as migracoes pendentes ate a versao mais recente
alembic upgrade head

# Reverter a ultima migracao em caso de anomalia durante deploy
alembic downgrade -1
\end{lstlisting}

\section{Runbook Operacional e Monitoramento}
\subsection{Monitoramento de Saúde e Alertas}
\begin{itemize}
    \item Endpoint de Healthcheck: \texttt{GET /api/v1/health} (verifica conectividade com o banco PostgreSQL e serviços de cache);
    \item Logs Estruturados: Saída em formato JSON com nível (\texttt{INFO}, \texttt{WARNING}, \texttt{ERROR}), timestamp UTC e \texttt{correlation\_id} para rastreabilidade de requisições distribuídas.
\end{itemize}

\subsection{Procedimentos de Contingência, Backup e Rollback}
\begin{table}[htbp]
\caption{Plano de Resposta a Incidentes Operacionais}
\label{tab:incidentes_runbook}
\centering
\small
\begin{tabularx}{\textwidth}{@{} l Y Y @{}}
\toprule
\textbf{Incidente} & \textbf{Ação Imediata} & \textbf{Procedimento de Recuperação} \\
\midrule
Falha pós-deploy & Reverter tráfego para a versão de imagem Docker anterior. & Acionar comando de rollback no orchestrator e reverter migração se necessário. \\
Degradação de banco & Identificar consultas lentas no \texttt{pg\_stat\_activity}. & Reiniciar pool de conexões e escalonar réplica de leitura. \\
Corrupção de dados & Isolar a aplicação em modo de manutenção. & Restaurar o backup incremental mais recente do Supabase/PostgreSQL. \\
\bottomrule
\end{tabularx}
\end{table}
